\documentclass[10pt,letterpaper]{article}
\usepackage[margin=0.8in]{geometry}
\usepackage[T1]{fontenc}
\usepackage{lmodern,microtype,enumitem}
\usepackage[hidelinks]{hyperref}
\newcommand{\claim}[2]{#2}
\setlength{\parindent}{0pt}
\setlength{\parskip}{5pt}
\setlist[itemize]{leftmargin=*,itemsep=2pt,topsep=3pt}
\begin{document}
\InputIfFileExists{notes-start.tex}{}{\setcounter{page}{5}}
\begin{center}
{\Large\bfseries Draft-paper submission notes}\\[3pt]
Learning a Writing Rubric from Edits: Context or Weights?
\end{center}

\textbf{Model paper.}
\claim{N01}{Amanda Bertsch et al., \emph{In-Context Learning with Long-Context Models: An In-Depth Exploration},
\href{https://arxiv.org/html/2405.00200v2}{arXiv:2405.00200v2} (2025).
We adapt its empirical comparison of contextual demonstrations and fine-tuning
over data budgets, retaining setup, comparative results, diagnostic analysis,
and a bounded conclusion. Our narrower question concerns editor corrections,
an inferred playbook, and validity in a tool-based writing task.}

\textbf{Adapted outline and visual roles.}
\begin{itemize}
\item \claim{N02}{\textbf{Abstract and introduction:} establish the feedback-storage question,
the common-review design, and the distinction between private rubric compliance
and public task validity. State the measured result with its scope.}
\item \claim{N03}{\textbf{Related work:} place inferred preferences, maintained playbooks,
low-rank adaptation, and combined representations in one compact comparison.
Move this discussion before the experimental design so readers understand why
the learned playbook is the primary comparator. The broader literature map is
kept outside the half-length paper.}
\item \claim{N04}{\textbf{Method:} describe the writing task, private rubric, shared reviews,
and context and weight adaptation. A workflow diagram shows the common archive
branching into a playbook and a LoRA adapter, then evaluation on fresh assignments.}
\item \claim{N05}{\textbf{Experiments:} describe fresh assignment splits, development
selection, explanatory controls, and paired inference. Expand validity and scope relative
to the model paper because the composite endpoint needs explicit denominators.}
\item \claim{N06}{\textbf{Results:} present all three review budgets and the prespecified
endpoint contrast. The central visual compares validity-adjusted rubric scores
and valid submission across budgets; showing both makes the tradeoff visible.
A compact endpoint table places explanatory controls beside the learned methods.
Text reports intervals and seed variation. A second table summarizes the separate
instruction probe.}
\item \claim{N07}{\textbf{Results, diagnostic comparisons:} use unedited and shuffled targets, a combined
adapter/playbook policy, and an explicit-rule privileged reference to challenge
the interpretation. These are exploratory follow-ups on the same test set,
not independent confirmations of a selected mechanism.}
\item \claim{N08}{\textbf{Discussion:} bound the conclusion by the synthetic review archive,
planted scorer, single model, three method seeds, and absence of human prose
ratings. Distinguish cumulative fresh-base adaptation from sequential learning.
Omit the model paper's attention and ordering studies because this experiment
does not measure those mechanisms; omit broad multi-model and multi-dataset
figures because our evidence is intentionally narrower.}
\item \claim{N09}{\textbf{References and audit:} cite the primary research sources in the
paper. Keep the detailed evidence, source locators, frozen implementation,
response transcripts, and claim-by-claim justifications in the companion files.}
\end{itemize}

\textbf{Revisions and status.}
\claim{N10}{The starting proposal at commit \texttt{db853b1} covered writing and Lean.
The study narrows the empirical question to learning a writing rubric, implements a
learned-playbook comparator, holds review availability fixed, selects
hyperparameters on development prompts, and separates the primary comparison
from diagnostic controls. A workflow diagram makes the two adaptation routes
explicit, and the results figure separates rubric scores from public validity.
The failed development validity safeguard remains explicit. The original
per-rule calibration goal was not met; harder-rule and all-rule outcomes
therefore supplement the composite score.
Lean transfer and optimizer-continuous learning are outside this draft's claims.
No instructor or classmate feedback was supplied, so none is represented as addressed.}

\textbf{Length and delivery.}
\claim{N11}{The user's half-length requirement overrides the ordinary full-paper length.
The measured body-word and page counts appear in \texttt{evidence/build-check.json}.
This combined PDF appends the model-paper analysis, outline, and revision notes
to the draft, following the
\href{https://web.stanford.edu/class/cs197/assignments/project.html\#draft-paper}{CS197 draft-paper requirements}.
The separate \texttt{CLAIMS.md} audit distinguishes measurements, implementation
facts, literature, inferences, and scope limits. The files are prepared for review;
they have not been uploaded to Overleaf or submitted to a class or venue.}
\end{document}
